Genomic prediction models are trained almost universally to minimize the mean squared error (MSE), yet they are evaluated, and used to make selection decisions, with the Pearson correlation $r$ between predicted and observed phenotypes and with the ranking of candidate genotypes. We do not claim novelty for the underlying identity: that maximizing $r$ is equivalent to minimizing the MSE of standardized variables, $\mathrm{MSE}(z_y,z_{\hat y})=2(1-r)$, is the elementary fact that $r$ is the inner product of $z$-scored vectors, and correlation-type, concordance (CCC), and differentiable ranking losses have been used before in machine learning and genomic prediction. Our contribution is instead a systematic, calibration-aware and selection-metric-aware empirical characterization of these losses for genomic prediction across $101$ trait--dataset combinations from $12$ species (EasyGeSe, CIMMYT wheat, SoyNAM), together with a closed-form post-hoc affine recalibration framing. We use the standardized-MSE identity only as a numerically stable training target (Pearson-loss implementation matched to a reference to single-precision tolerance, maximum deviation $5\times10^{-8}$ over $80$ random cases), and we note that the same scale- and offset-invariance lets a single post-hoc affine map restore magnitude with residual variance $\sigma_y^2(1-r^2)$ while leaving $r$ and all rankings unchanged whenever the fitted slope is positive (it is in $>99\%$ of folds; under $1\%$ show a sign flip when validation covariance is negative).

With architectures, splits and tuning held fixed so that only the loss differs, and comparing each architecture against its own MSE baseline, correlation-consistent and CCC losses raised predictive ability over MSE (pooled $\Delta r$ of $+0.038$ to $+0.044$; Holm $p<10^{-7}$ from a confirmatory mixed model) and selection-aware metrics (NDCG@$10$ $+0.016$; relative efficiency $+0.04$). The benefit is strongly architecture-dependent: large for the Transformer ($\Delta r=+0.10$, but estimated on only $49$ of the $101$ trait--datasets, from $4$ species) and CNN, and negligible for the MLP ($\Delta r=+0.003$, confidence interval including zero); the pooled $n=251$ is an unbalanced average over CNN ($101$), MLP ($101$) and Transformer ($49$) cells. Among the loss changes, only CCC reliably lowered affine-calibrated RMSE (pooled $-0.057$); the Pearson and hybrid losses left calibrated RMSE statistically null (confidence intervals spanning zero, despite a small sign-based Wilcoxon $p$). The gain weakens, and on some metrics reverses, under harder leave-family-out and cross-environment splits, but these contrasts rest on only $4$--$12$ non-independent cells and are underpowered rather than conclusive. No loss unseated the ridge/GBLUP baselines on mean rank. Because hyperparameters were selected on full-data validation under MSE, absolute accuracies may be mildly optimistic; this bias is symmetric within each architecture, so it does not affect the relative loss comparisons, and it works against, not for, the conclusion that neural networks do not outrank the linear baselines.